% ====================================================================
% ch3v22_bib.tex — Chapter 3(Si·혼합) 장별 참고문헌 (14종) — R1 자동 분할(cite 스캔 기반)
%   원장 = results/V1022_REFERENCE_LEDGER.md (V1 키만). 장 자기완결(리뷰 모음형) — 공통 키의 타 장 중복 수록은 의도(D22).
% ====================================================================
\begin{thebibliography}{99}
\bibitem{wen_huggins1981} C. J. Wen and R. A. Huggins, ``Chemical diffusion in intermediate phases in the lithium-silicon system,'' \emph{J. Solid State Chem.} \textbf{37}(3), 271--278 (1981). DOI: 10.1016/0022-4596(81)90487-4. [고온 평형 Li--Si 중간상 titration 원전.]
\bibitem{limthongkul2003} P. Limthongkul, Y.-I. Jang, N. J. Dudney, and Y.-M. Chiang, ``Electrochemically-driven solid-state amorphization in lithium-silicon alloys and implications for lithium storage,'' \emph{Acta Mater.} \textbf{51}(4), 1103 (2003). DOI: 10.1016/S1359-6454(02)00514-1.
\bibitem{li_dahn2007} J. Li and J. R. Dahn, ``An In Situ X-Ray Diffraction Study of the Reaction of Li with Crystalline Si,'' \emph{J. Electrochem. Soc.} \textbf{154}(3), A156--A161 (2007). DOI: 10.1149/1.2409862.
\bibitem{obrovac_christensen2004} M. N. Obrovac and L. Christensen, ``Structural Changes in Silicon Anodes during Lithium Insertion/Extraction,'' \emph{Electrochem. Solid-State Lett.} \textbf{7}(5), A93--A96 (2004). DOI: 10.1149/1.1652421. [Li$_{15}$Si$_4$ 준안정 결정화.]
\bibitem{chevrier_dahn2009} V. L. Chevrier and J. R. Dahn, ``First Principles Model of Amorphous Silicon Lithiation,'' \emph{J. Electrochem. Soc.} \textbf{156}(6), A454--A458 (2009). DOI: 10.1149/1.3111037. [비정질 경로의 경사 전위-조성 곡선.]
\bibitem{beaulieu2001} L. Y. Beaulieu, K. W. Eberman, R. L. Turner, L. J. Krause, and J. R. Dahn, ``Colossal Reversible Volume Changes in Lithium Alloys,'' \emph{Electrochem. Solid-State Lett.} \textbf{4}(9), A137 (2001). DOI: 10.1149/1.1388178.
\bibitem{sethuraman_stressevo2010} V. A. Sethuraman, M. J. Chon, M. Shimshak, V. Srinivasan, and P. R. Guduru, ``In situ measurements of stress evolution in silicon thin films during electrochemical lithiation and delithiation,'' \emph{J. Power Sources} \textbf{195}(15), 5062--5066 (2010). DOI: 10.1016/j.jpowsour.2010.02.013. [소성 유동 $\sim-1.75$ GPa.]
\bibitem{sethuraman_stresspot2010} V. A. Sethuraman, V. Srinivasan, A. F. Bower, and P. R. Guduru, ``In Situ Measurements of Stress-Potential Coupling in Lithiated Silicon,'' \emph{J. Electrochem. Soc.} \textbf{157}(11) (2010). DOI: 10.1149/1.3489378. [응력--전위 결합 $\sim100$--$120$ mV/GPa --- 쪽 범위는 사용 시 Crossref 최종 대조.]
\bibitem{liu_sizefracture2012} X. H. Liu, L. Zhong, S. Huang, S. X. Mao, T. Zhu, and J. Y. Huang, ``Size-Dependent Fracture of Silicon Nanoparticles During Lithiation,'' \emph{ACS Nano} \textbf{6}(2), 1522--1531 (2012). DOI: 10.1021/nn204476h. [임계 $\sim150$ nm.]
\bibitem{obrovac_chevrier2014} M. N. Obrovac and V. L. Chevrier, ``Alloy Negative Electrodes for Li-Ion Batteries,'' \emph{Chem. Rev.} \textbf{114}(23), 11444--11502 (2014). DOI: 10.1021/cr500207g. [합금 음극 총람 리뷰, tier B.]
\bibitem{verbrugge_lisi2016} M. W. Verbrugge, D. R. Baker, and X. Xiao, ``Formulation for the Treatment of Multiple Electrochemical Reactions and Associated Speciation for the Lithium-Silicon Electrode,'' \emph{J. Electrochem. Soc.} \textbf{163}(2), A262--A271 (2016). DOI: 10.1149/2.0581602jes. [다반응$\cdot$speciation --- 전하 보존 구조의 Si 실증(MSMR 계보).]
\bibitem{jiang_sihys2020} Y. Jiang, G. Offer, J. Jiang, M. Marinescu, and H. Wang, ``Voltage Hysteresis Model for Silicon Electrodes for Lithium Ion Batteries, Including Multi-Step Phase Transformations, Crystallization and Amorphization,'' \emph{J. Electrochem. Soc.} \textbf{167}(13), 130533 (2020). DOI: 10.1149/1945-7111/abbbba.
\bibitem{larchecahn1973} F. Larch\'e and J. W. Cahn, ``A linear theory of thermochemical equilibrium of solids under stress,'' \emph{Acta Metall.} \textbf{21}(8), 1051--1063 (1973). DOI: 10.1016/0001-6160(73)90021-7. [응력--조성 결합 화학퍼텐셜의 이론 원전.]
\bibitem{koebbing2024} L. K\"obbing, A. Latz, and B. Horstmann, ``Voltage Hysteresis of Silicon Nanoparticles: Chemo-Mechanical Particle-SEI Model,'' \emph{Adv. Funct. Mater.}, 2308818. DOI: 10.1002/adfm.202308818. [SEI 점탄소성 소산 히스 모델 --- 인쇄 권$\cdot$호는 사용 시 Crossref 최종 대조.]

% ---- R5 신규 등재(19종 — comp_R4 검증·Crossref 재생성·CHERRYPICK_R5) ----
\bibitem{wang_asi2013} J. W. Wang, Y. He, F. Fan, X. H. Liu, S. Xia, Y. Liu, C. T. Harris, H. Li, J. Y. Huang, S. X. Mao, T. Zhu, ``Two-Phase Electrochemical Lithiation in Amorphous Silicon,'' \emph{Nano Lett.} \textbf{13}, 709--715 (2013). DOI: 10.1021/nl304379k.
\bibitem{mcdowell_coreshell2013} M. T. McDowell, S. W. Lee, J. T. Harris, B. A. Korgel, C. Wang, W. D. Nix, Y. Cui, ``In Situ TEM of Two-Phase Lithiation of Amorphous Silicon Nanospheres,'' \emph{Nano Lett.} \textbf{13}, 758--764 (2013). DOI: 10.1021/nl3044508.
\bibitem{ogata_nmr2014} K. Ogata, E. Salager, C. J. Kerr, A. E. Fraser, C. Ducati, A. J. Morris, S. Hofmann, C. P. Grey, ``Revealing lithium–silicide phase transformations in nano-structured silicon-based lithium ion batteries via in situ NMR spectroscopy,'' \emph{Nat. Commun.} \textbf{5}, 3217 (2014). DOI: 10.1038/ncomms4217.
\bibitem{miyachi_sio2005} M. Miyachi, H. Yamamoto, H. Kawai, T. Ohta, M. Shirakata, ``Analysis of SiO Anodes for Lithium-Ion Batteries,'' \emph{J. Electrochem. Soc.} \textbf{152}, A2089 (2005). DOI: 10.1149/1.2013210.
\bibitem{kitada_sio2019} K. Kitada, O. Pecher, P. C. M. M. Magusin, M. F. Groh, R. S. Weatherup, C. P. Grey, ``Unraveling the Reaction Mechanisms of SiO Anodes for Li-Ion Batteries by Combining in Situ $^{7}$Li and ex Situ $^{7}$Li/$^{29}$Si Solid-State NMR Spectroscopy,'' \emph{J. Am. Chem. Soc.} \textbf{141}, 7014--7027 (2019). DOI: 10.1021/jacs.9b01589.
\bibitem{zhang_sio2018} L. Zhang, Y. Qin, Y. Liu, Q. Liu, Y. Ren, A. N. Jansen, W. Lu, ``Capacity Fading Mechanism and Improvement of Cycling Stability of the SiO Anode for Lithium-Ion Batteries,'' \emph{J. Electrochem. Soc.} \textbf{165}, A2102--A2107 (2018). DOI: 10.1149/2.0431810jes.
\bibitem{yom_sio2016} J. H. Yom, S. W. Hwang, S. M. Cho, W. Y. Yoon, ``Improvement of irreversible behavior of SiO anodes for lithium ion batteries by a solid state reaction at high temperature,'' \emph{J. Power Sources} \textbf{311}, 159--166 (2016). DOI: 10.1016/j.jpowsour.2016.02.025.
\bibitem{andersen_sic2019} H. F. Andersen, C. E. L. Foss, J. Voje, R. Tronstad, T. Mokkelbost, P. E. Vullum, A. Ulvestad, M. Kirkengen, J. P. Mæhlen, ``Silicon-Carbon composite anodes from industrial battery grade silicon,'' \emph{Sci. Rep.} \textbf{9}, 14814 (2019). DOI: 10.1038/s41598-019-51324-4.
\bibitem{naboka_sic2021} O. Naboka, C. Yim, Y. Abu-Lebdeh, ``Practical Approach to Enhance Compatibility in Silicon/Graphite Composites to Enable High-Capacity Li-Ion Battery Anodes,'' \emph{ACS Omega} \textbf{6}, 2644--2654 (2021). DOI: 10.1021/acsomega.0c04811.
\bibitem{lee_sic2025} Y. Lee, S. C. Hong, K. Kim, J. Park, J. Kim, B. Kang, Y. Kim, C. Lee, S. Son, B. Cho, J. Kim, S. Woo, G. Lee, K. Kang, ``Multiscale Carbon‐Integrated Silicon Anode for Stable Cycling Under Practical Lithium‐Ion Battery Conditions,'' \emph{Adv. Energy Mater.} e04250 (2025). DOI: 10.1002/aenm.202504250.
\bibitem{bohm_entropy2024} K. Böhm, S. Zintel, P. Ganninger, J. Jäger, T. Markus, D. Henriques, ``Exploring the Impact of State of Charge and Aging on the Entropy Coefficient of Silicon–Carbon Anodes,'' \emph{Energies} \textbf{17}, 5790 (2024). DOI: 10.3390/en17225790.
\bibitem{arnot_calorimetry2021} D. J. Arnot, E. Allcorn, K. L. Harrison, ``Effect of Temperature and FEC on Silicon Anode Heat Generation Measured by Isothermal Microcalorimetry,'' \emph{J. Electrochem. Soc.} \textbf{168}, 110509 (2021). DOI: 10.1149/1945-7111/ac315c.
\bibitem{bohm_thermal2025} K. Böhm, C. Bracht, C. Kallfaß, T. Markus, D. Henriques, ``Temperature-Dependent Thermal Effects and Capacity Characteristics of Silicon Anodes in Lithium-Ion Batteries,'' \emph{J. Electrochem. Soc.} \textbf{172}, 050537 (2025). DOI: 10.1149/1945-7111/adda7a.
\bibitem{wojtala_entropy2022} M. E. Wojtala, A. A. Zülke, R. Burrell, M. Nagarathinam, G. Li, H. E. Hoster, D. A. Howey, M. P. Mercer, ``Entropy Profiling for the Diagnosis of NCA/Gr-SiOx Li-Ion Battery Health,'' \emph{J. Electrochem. Soc.} \textbf{169}, 100527 (2022). DOI: 10.1149/1945-7111/ac87d1.
\bibitem{gautam_blend2024} M. Gautam, G. K. Mishra, K. Bhawana, C. S. Kalwar, D. Dwivedi, A. Yadav, S. Mitra, ``Relationship between Silicon Percentage in Graphite Anode to Achieve High-Energy-Density Lithium-Ion Batteries,'' \emph{ACS Appl. Mater. Interfaces} \textbf{16}, 45809--45820 (2024). DOI: 10.1021/acsami.4c10178.
\bibitem{ai_composite2022} W. Ai, N. Kirkaldy, Y. Jiang, G. Offer, H. Wang, B. Wu, ``A composite electrode model for lithium-ion batteries with silicon/graphite negative electrodes,'' \emph{J. Power Sources} \textbf{527}, 231142 (2022). DOI: 10.1016/j.jpowsour.2022.231142.
\bibitem{chatzogiannakis_blend2025} D. Chatzogiannakis, I. Ghilescu, G. Giannadaki, M. Cabello, M. Casas‐Cabanas, M. R. Palacin, ``Decoupling Silicon and Graphite Contribution in High‐Silicon Content Composite Electrodes,'' \emph{Batteries \& Supercaps} \textbf{8}, e202500104 (2025). DOI: 10.1002/batt.202500104.
\bibitem{moyassari_blend2022} E. Moyassari, T. Roth, S. K\"ucher, C.-C. Chang, S.-C. Hou, F. B. Spingler, and A. Jossen, ``The Role of Silicon in Silicon-Graphite Composite Electrodes Regarding Specific Capacity, Cycle Stability, and Expansion,'' \emph{J. Electrochem. Soc.} \textbf{169}(1), 010504 (2022). DOI: 10.1149/1945-7111/ac4545. [Si 0--20 wt\% 스윕 딜라토메트리$\cdot$전기화학 동시 실측.]
\bibitem{zhan_siox2026} M. Zhan, B. Jin, N. Stapf, J. Meyer, K. P. Birke, A. Fill, ``Unraveling the vertical expansion and hysteresis in SiOx/graphite composite electrodes via in-situ dilatometry and cracking origins via X-ray diffraction and scanning electron microscopy,'' \emph{J. Energy Storage} \textbf{154}, 121227 (2026). DOI: 10.1016/j.est.2026.121227.
\bibitem{tu_blend2024} M. Tu, T. Dao, M. W. Verbrugge, B. Koch, ``Mathematical Model for a Lithium-Ion Battery with a SiO/Graphite Blended Electrode Based on a Reduced Order Model Derived Using Perturbation Theory,'' \emph{J. Electrochem. Soc.} \textbf{171}, 050539 (2024). DOI: 10.1149/1945-7111/ad4823.
\bibitem{artrith2018} N. Artrith, A. Urban, and G. Ceder, ``Constructing first-principles phase diagrams of amorphous Li$_x$Si using machine-learning-assisted sampling with an evolutionary algorithm,'' \emph{J. Chem. Phys.} \textbf{148}(24), 241711 (2018). DOI: 10.1063/1.5017661. [비정질 Li$_x$Si 제일원리/ML 상도표 --- 연속 경사 $U(x)$ 독립 재현(단일상 고용체 근거). tier B.]
\bibitem{verbrugge2017} M. W. Verbrugge, D. R. Baker, B. J. Koch, X. Xiao, and W. Gu, ``Thermodynamic Model for Substitutional Materials: Application to Lithiated Graphite, Spinel Manganese Oxide, Silicon, and Their Alloys,'' \emph{J. Electrochem. Soc.} \textbf{164}(11), E3243--E3253 (2017). DOI: 10.1149/2.0341708jes. [substitutional 열역학 $\omega$(상호작용)-형식 --- Frumkin 커널의 정칙용액 골격 계보(MSMR $\omega$ 형식 원전). tier B.]
\end{thebibliography}
